% part: intuitionistic-logic
% chapter: introduction
% section: natural-deduction

\documentclass[../../../include/open-logic-section]{subfiles}

\begin{document}

\olfileid{int}{int}{ntd}

\olsection{طبيعي استنتاج}

له $\FalseCl$ قاعدو پرته طبيعي استنتاج د شهودي منطق يو معياري
!!{derivation} نظام دے۔ دلته قاعدې بيا راوړو او د BHK تفسير له مخې
ئې انګېزه څرګندوو۔ په هر حالت کښې قاعده داسې ګڼلے شو چې دا پايله
راوباسي: که مقدمې جوړونې ولري، نو نتيجه ئې هم جوړونه لري۔

څرنګه چې د طبيعي استنتاج په !!{derivation}s کښې نۀ ختم شوي
فرضونه وي، نو د بېلګې په توګه له !!{undischarged} فرضونو~$\Gamma$
څخه د $!A$ داسې !!a{derivation} بايد يو تابع وګڼو چې د هر
$!B \in \Gamma$ جوړونې د~$!A$ په جوړونه اړوي۔ که د~$!A$
داسې !!a{derivation} وي چې هېڅ !!{undischarged} فرض ونۀ لري،
نو د BHK تفسير په معنا د~$!A$ جوړونه شته۔ خو د بحث د اسانۍ لپاره
به چې اړتيا نۀ وي، $\Gamma$ پرېږدو۔

يو فرض~$!A$ پخپله له !!{undischarged} فرض~$!A$ څخه د~$!A$
!!a{derivation} دے۔ دا له BHK تفسير سره سمون لري: د جوړونو عيني
تابع د~$!A$ هره جوړونه د~$!A$ په جوړونه اړوي۔

\subsection{عطف}

\begin{defish}
\AxiomC{$!A$}
\AxiomC{$!B$}
\RightLabel{\Intro{\land}}
\BinaryInfC{$!A \land !B$}
\DisplayProof
\hfill
\begin{tabular}{l}
\AxiomC{$!A \land !B$}
\RightLabel{\Elim{\land}}
\UnaryInfC{$!A$}
\DisplayProof
\\[3ex]
\AxiomC{$!A \land !B$}
\RightLabel{\Elim{\land}}
\UnaryInfC{$!B$}
\DisplayProof
\end{tabular}
\end{defish}

فرض کړئ چې په ترتيب د $!A_1$ او $!A_2$ جوړونې $N_1$ او $N_2$
لرو۔ نو د $!A_1 \land !A_2$ جوړونه هم لرو، يعنې جوړه
$\tuple{N_1, N_2}$۔

د BHK تفسير له مخې د $!A_1 \land !A_2$ جوړونه يوه جوړه
$\tuple{N_1, N_2}$ ده۔ نو فرض کړئ چې داسې جوړه لرو۔ بيا د عطف
د دواړو غړو جوړونې هم لرو: $N_1$ د~$!A_1$ او $N_2$ د~$!A_2$
جوړونه ده۔ د ژباړې د نسخې څرګندونه (OLINT-003): په منجمده سرچينه
کښې د همدې جملې دويم عطف شوی غړی په تېروتنه د لومړي غړي په نښه
بيا ليکل شوی؛ د جوړې د دويمې برخې او د ورپسې جملې له مخې دلته
دويم غړی مراد دے۔

\subsection{شرط}

\begin{defish}
  \AxiomC{$\Discharge{!A}{u}$}
  \DeduceC{$!B$}
  \DischargeRule{$\Intro{\lif}$}{u}
  \UnaryInfC{$!A \lif !B$}
  \DisplayProof
\hfill
  \AxiomC{$!A \lif !B$}
  \AxiomC{$!A$}
  \RightLabel{$\Elim{\lif}$}
  \BinaryInfC{$!B$}
  \DisplayProof
\end{defish}

که له !!{undischarged} فرض~$!A$ څخه د~$!B$ !!a{derivation}
ولرو، نو داسې تابع~$f$ شته چې د~$!A$ جوړونې د~$!B$ په جوړونو
اړوي۔ همدا تابع د $!A \lif !B$ جوړونه ده۔ نو که د $\Intro\lif$
مقدمې ته د~$!A$ پر جوړونه مشروطه جوړونه وي، نتيجه $!A
\lif !B$ هم جوړونه لري۔

بل خوا فرض کړئ چې $N$ د $!A$ جوړونه او $f$ د $!A \lif !B$
جوړونه ده۔ د $!A \lif !B$ جوړونه داسې تابع ده چې د $!A$
جوړونې د~$!B$ په جوړونو اړوي۔ نو $f(N)$ د~$!B$ جوړونه ده؛
يعنې د $\Elim\lif$ نتيجه جوړونه لري۔

\subsection{فصل}

\begin{defish}
\begin{tabular}{l}
\AxiomC{$!A$}
\RightLabel{\Intro{\lor}}
\UnaryInfC{$!A \lor !B$}
\DisplayProof
\\[3ex]
\AxiomC{$!B$}
\RightLabel{\Intro{\lor}}
\UnaryInfC{$!A \lor !B$}
\DisplayProof
\end{tabular}
\hfill
\AxiomC{$!A \lor !B$}
\AxiomC{$\Discharge{!A}{n}$}
\DeduceC{$!C$}
\AxiomC{$\Discharge{!B}{n}$}
\DeduceC{$!C$}
\DischargeRule{\Elim{\lor}}{n}
\TrinaryInfC{$!C$}
\DisplayProof
\end{defish}

که د~$!A_i$ جوړونه~$N_i$ ولرو، هغې ته د~$!A_1 \lor !A_2$
جوړونې~$\tuple{i, N_i}$ بڼه ورکولے شو۔ بل خوا فرض کړئ چې د
~$!A_1 \lor !A_2$ جوړونه لرو، يعنې يوه جوړه~$\tuple{i, N_i}$
چې $N_i$ د~$!A_i$ جوړونه ده۔ همدارنګه تابعې $f_1$ او $f_2$
لرو چې په ترتيب د $!A_1$ او $!A_2$ جوړونې د~$!C$ په جوړونو
اړوي۔ نو $f_i(N_i)$ د~$!C$ جوړونه، يعنې د~$\Elim\lor$ نتيجه ده۔

\subsection{محال}

\begin{defish}
  \AxiomC{$\lfalse$}
  \RightLabel{\FalseInt}
  \UnaryInfC{$!A$}
  \DisplayProof
\end{defish}

که له !!{undischarged} فرضونو~$!B_1$، \dots، $!B_n$ څخه د
~$\lfalse$ !!a{derivation} ولرو، نو تابع~$f(M_1,
\dots, M_n)$ د $!B_1$، \dots، $!B_n$ جوړونې د~$\lfalse$
په جوړونه اړوي۔ څرنګه چې $\lfalse$ هېڅ جوړونه نۀ لري، نو د
$!B_1$، \dots، $!B_n$ ټولو جوړونې هم يوځاے نۀ شي موجودېدے۔
له دې امله $f$ دا خاصيت هم لري چې \emph{که} $M_1$، \dots،
$M_n$ په ترتيب د $!B_1$، \dots، $!B_n$ جوړونې وي،
\emph{نو} $f(M_1, \dots, M_n)$ د~$!A$ جوړونه ده۔

\subsection{د $\lnot$ قاعدې}

ځکه چې $\lnot !A$ د $!A \lif \lfalse$ په توګه تعريفېږي، په
دقيقه معنا د~$\lnot$ قاعدو ته اړتيا نۀ لرو۔ خو که ئې وغواړو،
بڼه به ئې داسې وي:

\begin{defish}
\AxiomC{$\Discharge{!A}{n}$}
\noLine
\DeduceC{$\lfalse$}
\DischargeRule{\Intro{\lnot}}{n}
\UnaryInfC{$\lnot !A$}
\DisplayProof
\hfill
\AxiomC{$\lnot !A$}
\AxiomC{$!A$}
\RightLabel{\Elim{\lnot}}
\BinaryInfC{$\lfalse$}
\DisplayProof
\end{defish}


\subsection{د \usetoken{P}{derivation} بېلګې}

\begin{enumerate}
\item $\Proves !A \lif (\lnot !A \lif \lfalse)$، يعنې $\Proves !A
  \lif ((!A \lif \lfalse) \lif \lfalse)$
  \begin{prooftree}
    \AxiomC{$\Discharge{!A}{2}$}
    \AxiomC{$\Discharge{!A \lif \lfalse}{1}$}
    \RightLabel{\Elim\lif}
    \BinaryInfC{$\lfalse$}
    \DischargeRule{\Intro\lif}{1}
    \UnaryInfC{$(!A \lif \lfalse)\lif \lfalse$}
    \DischargeRule{\Intro\lif}{2}
    \UnaryInfC{$!A \lif (!A \lif \lfalse) \lif \lfalse$}
  \end{prooftree}

\item $\Proves ((!A \land !B) \lif !C) \lif (!A \lif (!B \lif !C))$
  \begin{prooftree}
    \AxiomC{$\Discharge{(!A \land !B) \lif !C}{3}$}
    \AxiomC{$\Discharge{!A}{2}$}
    \AxiomC{$\Discharge{!B}{1}$}
    \RightLabel{\Intro\land}
    \BinaryInfC{$!A \land !B$}
    \RightLabel{\Elim\lif}
    \BinaryInfC{$!C$}
    \DischargeRule{\Intro\lif}{1}
    \UnaryInfC{$!B \lif !C$}
    \DischargeRule{\Intro\lif}{2}
    \UnaryInfC{$!A \lif (!B \lif !C)$}
    \DischargeRule{\Intro\lif}{3}
    \UnaryInfC{$((!A \land !B) \lif !C) \lif (!A \lif (!B \lif !C))$}
  \end{prooftree}

\item $\Proves \lnot(!A \land \lnot !A)$، يعنې $\Proves (!A \land (!A
  \lif \lfalse))\lif \lfalse$
  \begin{prooftree}
    \AxiomC{$\Discharge{!A \land (!A \lif \lfalse)}{1}$}
    \RightLabel{\Elim\land}
    \UnaryInfC{$!A \lif \lfalse$}
    \AxiomC{$\Discharge{!A \land (!A \lif \lfalse)}{1}$}
    \RightLabel{\Elim\land}
    \UnaryInfC{$!A$}
    \RightLabel{\Elim\lif}
    \BinaryInfC{$\lfalse$}
    \DischargeRule{\Intro\lif}{1}
    \UnaryInfC{$(!A \land (!A \lif \lfalse))\lif \lfalse$}
  \end{prooftree}

\item $\Proves \lnot\lnot(!A \lor \lnot !A)$، يعنې $\Proves ((!A \lor
  (!A \lif \lfalse))\lif \lfalse)\lif \lfalse$
  \begin{prooftree}\hskip-3em
    \AxiomC{$\Discharge{(!A \lor (!A \lif \lfalse))\lif \lfalse}{2}$}
    \AxiomC{$\Discharge{(!A \lor (!A \lif \lfalse))\lif \lfalse}{2}$}
    \AxiomC{$\Discharge{!A}{1}$}
    \RightLabel{\Intro\lor}
    \UnaryInfC{$!A \lor (!A \lif \lfalse)$}
    \RightLabel{\Elim\lif}
    \BinaryInfC{$\lfalse$}
    \DischargeRule{\Intro\lif}{1}
    \UnaryInfC{$!A \lif \lfalse$}
    \RightLabel{\Intro\lor}
    \UnaryInfC{$!A \lor (!A \lif \lfalse)$}
    \RightLabel{\Elim\lif}
    \insertBetweenHyps{\hskip -4em}
    \BinaryInfC{$\lfalse$}
    \DischargeRule{\Intro\lif}{2}
    \UnaryInfC{$((!A \lor (!A \lif \lfalse))\lif \lfalse)\lif \lfalse$}
  \end{prooftree}
\end{enumerate}

\begin{prop}
  که په شهودي منطق کښې د طبيعي استنتاج له مخې $\Gamma \Proves !A$
  وي، نو په کلاسيک منطق کښې هم $\Gamma \Proves !A$ وي۔ په تېره بيا،
  که $!A$ شهودي قضيه وي، نو کلاسيکه قضيه هم ده۔
\end{prop}

\begin{proof}
  د طبيعي استنتاج هره قاعده په کلاسيک طبيعي استنتاج کښې هم قاعده ده؛
  نو د شهودي منطق هر !!{derivation} په کلاسيک منطق کښې هم
  !!a{derivation} دے۔
\end{proof}

\begin{prob}
  په شهودي منطق کښې د لاندې !!{formula}s !!{derivation}s ورکړئ:
  \begin{enumerate}
    \item $(\lnot !A \lor !B) \lif (!A \lif !B)$
    \item $\lnot\lnot\lnot !A \lif \lnot !A$
    \item $\lnot\lnot (!A \land !B) \liff (\lnot\lnot !A\land \lnot \lnot !B)$
    \item $\lnot (!A \lor !B) \liff (\lnot !A \land \lnot !B)$
    \item $(\lnot !A \lor \lnot !B) \lif \lnot (!A \land !B)$
    \item $\lnot\lnot ( !A \land !B) \lif  (\lnot\lnot !A \lor
    \lnot\lnot !B)$
  \end{enumerate}
\end{prob}

\end{document}
